\documentclass[11pt,a4paper]{article}
\usepackage[utf8]{inputenc}
\usepackage{amsmath,amssymb}
\usepackage{geometry}
\usepackage{booktabs}
\usepackage{xcolor}
\usepackage{listings}
\usepackage{hyperref}

\geometry{margin=1in}

\title{\textbf{Day 2: Machine Learning and Neural Networks}}
\author{Data Science Foundation Series}
\date{\today}

\begin{document}

\maketitle

\section{AI vs. Machine Learning vs. Deep Learning}
\begin{itemize}
    \item \textbf{Artificial Intelligence (AI):} Any computational system mimicking human decision-making. Can be rule-based without learning.
    \item \textbf{Machine Learning (ML):} The computer discovers predictive statistical rules directly from past empirical data.
    \item \textbf{Deep Learning (DL):} Multi-layered artificial neural networks capable of automatic feature extraction from unstructured data (images, text, audio).
\end{itemize}

\section{Supervised vs. Unsupervised Learning}
\begin{itemize}
    \item \textbf{Supervised Learning:} The dataset contains ground truth target labels $y$. The goal is to learn a mapping $f(X) \approx y$.
    \item \textbf{Unsupervised Learning:} No target labels exist. The goal is discovering latent clusters or dimensionality reductions.
\end{itemize}

\section{Regression vs. Classification}
\begin{itemize}
    \item \textbf{Regression:} Predicting continuous real numbers $y \in \mathbb{R}$ (e.g. Price, Temperature).
    \item \textbf{Classification:} Predicting discrete category labels $y \in \{0, 1, \dots, K-1\}$.
\end{itemize}

\section{The Five Core Algorithms}
\begin{enumerate}
    \item \textbf{Linear Regression:} Fits the hyper-plane minimizing sum of squared residuals: $y = \mathbf{w}^T\mathbf{x} + b$.
    \item \textbf{Logistic Regression:} Models binary class probability via Sigmoid link function: $P(y=1|\mathbf{x}) = \frac{1}{1 + e^{-(\mathbf{w}^T\mathbf{x} + b)}}$.
    \item \textbf{Decision Tree:} Hierarchical recursive splitting based on impurity criteria (Gini or Entropy).
    \item \textbf{Random Forest:} Bagging ensemble of de-correlated decision trees with feature sub-sampling.
    \item \textbf{K-Means:} Partitioning algorithm iteratively minimizing within-cluster sum of squares to $K$ centroids.
\end{enumerate}

\section{Model Evaluation \& Overfitting}
\subsection{Regression Metrics}
\begin{equation}
\text{MAE} = \frac{1}{N}\sum_{i=1}^N |y_i - \hat{y}_i|, \quad \text{RMSE} = \sqrt{\frac{1}{N}\sum_{i=1}^N (y_i - \hat{y}_i)^2}, \quad R^2 = 1 - \frac{\sum (y_i - \hat{y}_i)^2}{\sum (y_i - \bar{y})^2}
\end{equation}

\subsection{Classification Metrics}
\begin{equation}
\text{Precision} = \frac{TP}{TP + FP}, \quad \text{Recall} = \frac{TP}{TP + FN}, \quad F_1 = 2 \times \frac{\text{Precision} \times \text{Recall}}{\text{Precision} + \text{Recall}}
\end{equation}

\section{Neural Networks \& Activation Functions}
\subsection{The Artificial Neuron}
\begin{equation}
z = \sum_{j=1}^d w_j x_j + b, \quad a = \sigma(z)
\end{equation}

\subsection{Activation Functions}
\begin{itemize}
    \item \textbf{ReLU:} $f(x) = \max(0, x)$. Provides non-linear piecewise representation avoiding vanishing gradients.
    \item \textbf{Sigmoid:} $\sigma(x) = \frac{1}{1 + e^{-x}}$.
\end{itemize}

\subsection{Gradient Descent}
Weight updates in the direction of the negative gradient:
\begin{equation}
\mathbf{w}^{(t+1)} = \mathbf{w}^{(t)} - \eta \nabla_{\mathbf{w}} \mathcal{L}(\mathbf{w})
\end{equation}
where $\eta$ is the learning rate.

\end{document}
